\chapter{Závěr}
\label{8-zaver}

V~rámci této práce byly analyzovány vybrané výsledky zpracování metody \zkratka{DORIS} analytického centra na Geodetické observatoři Pecný VÚGTK. Hodnoceny byly dva typy produktů, a to časové řady souřadnic stanic vysílačů a přesné dráhy družic vybavených přijímači. U~časových řad souřadnic stanic byly za období 2015--2025 posuzovány všechny stanice, u~nichž bylo indikováno podezřelé chování, přičemž pozornost byla věnována zejména dlouhodobým trendům, diskontinuitám a periodickým složkám. Pro analýzu přesných drah družic byly zvoleny družice cs2, h2c, h2d, ja3, s3a, s3b, s6a a srl. Analyzováno bylo prvních třicet dní roku 2024. Cílem bylo porovnat řešení poskytované Geodetickou observatoří Pecný s~řešením \zk{SSA}. Porovnání bylo provedeno dvěma odlišnými přístupy, a to pomocí dynamické parametrizace dráhy v~\zkratka{BSW} a pomocí Hermitovy interpolace souřadnic v~daných epochách. Součástí práce bylo rovněž vytvoření knihovny v~jazyce Python, která umožňuje jednotné zpracování dat systému \zk{DORIS}, jejich analýzu a generování výstupů.

Analýza časových řad souřadnic stanic umožnila posoudit stabilitu vybraných stanic systému \zkratka{DORIS} a identifikovat jejich dlouhodobé trendy, diskontinuity a periodické složky. U~většiny sledovaných stanic bylo indikované podezřelé chování potvrzeno, pravděpodobnou příčinu se však podařilo určit pouze v~některých případech. Nebyly zjištěny zásadní nekonzistence, které by ukazovaly na obecný problém řešení \zk{GOP}. Detekované odchylky mají převážně lokální charakter a jejich interpretace závisí na konkrétní stanici, období a porovnání s~dalšími dostupnými řešeními.

V~části věnované drahám družic bylo provedeno porovnání řešení \zk{GOP} a \zk{SSA} v~lokálním systému \zkratka{RTN}. Výsledky ukázaly, že rozdíly mezi řešeními se většinou pohybují na úrovni centimetrů, přičemž v~radiálním směru, který je u~altimetrických družic často nejdůležitější, dosahují typicky vyšších milimetrů až přibližně 1~cm. Radiální složka současně vykazovala u~všech družic nejnižší a nejstabilnější hodnoty \zk{RMS}, zatímco normálová a transverzální složka dosahovaly vyšších hodnot s~výraznější variabilitou. Ve spektrálních analýzách rozdílů se u~většiny družic projevila periodická složka odpovídající přibližně oběžné době. Oba použité přístupy k~porovnání drah -- Hermitova interpolace a dynamická parametrizace v~\zk{BSW} -- poskytly srovnatelné výsledky se~vzájemnými odchylkami na úrovni jednotek milimetrů, čímž se potvrdila použitelnost interpolační metody.

Výsledky práce přispívají k~hodnocení kvality řešení poskytovaného analytickým centrem na Geodetické observatoři Pecný. Jsou shrnuty výsledky analýzy stability vybraných časových řad souřadnic stanic systému \zkratka{DORIS} a~současně porovnány přesné dráhy družic mezi řešeními \zk{GOP} a~\zk{SSA}. Získané poznatky mají význam pro posouzení spolehlivosti výstupů \zk{GOP}, které jsou využívány mimo jiné při realizaci globálního referenčního rámce (\zk{ITRF}). Vytvořená knihovna umožňuje systema\-tickou analýzu těchto produktů díky jednotnému zpracování dat, porovnání výsledků jednotlivých řešení a~automatizovanému generování výstupů.

Práce je zatížena několika omezeními, která vyplývají jak z~charakteru vstupních dat, tak z~použitých metod. V~případě analýzy časových řad neposkytují některé stanice data za celé analyzované období. U~většiny případů detekovaného nestandardního chování se nepodařilo s~jistotou určit jeho příčinu. Analýza přesných drah družic byla provedena na omezeném časovém úseku, což neumožňuje plně vyhodnotit dlouhodobé chování. Porovnání drah bylo provedeno mezi dvěma nezávislými řešeními. Nalezené rozdíly proto nelze jednoznačně přiřadit ani jednomu z~nich.

Práci je možné dále rozšířit o~analýzu většího počtu stanic a~družic, delších časových období a~dalších řešení analytických center. Vytvořenou knihovnu lze zároveň doplnit o~další metody analýzy. Do budoucna by bylo možné nad knihovnou vytvořit grafické uživatelské rozhraní pro jednodušší analýzu dat.